\begin{figure}[ht]
    \centering
    \resizebox{0.8\textwidth}{!}{
        \begin{tikzpicture}[
            >={Stealth[length=3mm]},
            actor/.style={
                    ellipse, draw=blue!60, fill=blue!8,
                    minimum width=3cm, minimum height=1cm,
                    align=center, font=\small
                },
            process/.style={
                    rectangle, rounded corners=4pt, draw, fill=orange!10,
                    minimum width=3cm, minimum height=1cm,
                    align=center, font=\small
                },
            datastore/.style={
                    rectangle, draw, dashed, fill=purple!10, rounded corners=2pt,
                    minimum width=3cm, minimum height=1cm,
                    align=center, font=\small\itshape
                },
            output/.style={
                    rectangle, rounded corners=4pt, draw=green!50!black, fill=green!8,
                    minimum width=3cm, minimum height=1cm,
                    align=center, font=\small
                },
            label/.style={
                    font=\scriptsize, fill=white, inner sep=1pt
                },
            ]

            % Actor
            \node[actor] (learner) {Learner};

            % Row 1: Input processing
            \node[process, below=of learner] (intent) {Input Analysis};

            % Row 2: Two paths
            \node[process, below left=and 2.5cm of intent] (rag) {RAG Engine};
            \node[process, below right=and 2.5cm of intent] (difficulty) {Difficulty\\Matching};

            % Row 3: Generation / Reward
            \node[process, below=of rag] (llm) {MiniGuide LLM};
            \node[process, below=of difficulty] (reward) {Card Activation};

            % Row 4: Outputs
            \node[output, below=of llm] (question) {Guiding Question};
            \node[output, below=of reward] (gamestate) {Updated Game State};

            % Data stores
            \node[datastore, below=of intent] (curriculum) {Curriculum \&\\Heuristics KB};
            \node[datastore, below=of curriculum] (usermodel) {Dynamic User Model};

            % Arrows
            \draw[->] (learner) -- node[label, right] {Submit Answer / Ask for Help} (intent);
            \draw[->] (intent) -- node[label, above left] {Needs Help} (rag);
            \draw[->] (intent) -- node[label, above right] {Answer Submitted} (difficulty);

            % Left path: Help
            \draw[->] (rag) -- (llm);
            \draw[->] (llm) -- (question);
            \draw[->] (curriculum) -- (rag);

            % Right path: Game
            \draw[->] (difficulty) -- (reward);
            \draw[->] (reward) -- (gamestate);
            \draw[->] (curriculum) -- (difficulty);

            % Outputs back to learner
            \draw[->] (question.south) -- ++(0,-0.5) -| ([xshift=-7cm]learner.west) -- (learner.west);
            \draw[->] (gamestate.south) -- ++(0,-0.5) -| ([xshift=7cm]learner.east) -- (learner.east);

            % User model connections
            \draw[->, dashed] (usermodel) -- node[label, above] {ZPD} (llm);
            \draw[->, dashed] (usermodel) -- node[label, above left] {ZPD} (difficulty);
            \draw[->, dashed] (question.east) -- ++(1.5,0) node[label, above] {Update ZPD} -| ([xshift=-0.75cm]usermodel.south);
            \draw[->, dashed] (gamestate.west) -- ++(-1.5,0) node[label, above] {Update ZPD} -| ([xshift=0.75cm]usermodel.south);

        \end{tikzpicture}
    }

    \caption{Data flow diagram of the system architecture. Solid arrows indicate the primary data flow; dashed arrows indicate reads from and writes to the Dynamic User Model.}
    \label{fig:data_flow_diagram}
\end{figure}
